\documentclass[11pt]{article}
\usepackage[utf8]{inputenc}
\usepackage[T1]{fontenc}
\usepackage{newtxtext,newtxmath}
\usepackage{amsmath}
\usepackage{multicol}
\usepackage{geometry}
\usepackage{xcolor}
\usepackage{titlesec}

\geometry{a4paper, left=1cm, right=1cm, top=0.5cm, bottom=1.2cm}
\setlength{\columnseprule}{0.4pt}

\titleformat{\section}{\normalfont\Large\color{blue}}{\thesection}{1em}{}

\title{\textcolor{blue}{Questionário com Respostas \\ \textit{Nosedive} (Black Mirror)}}
\author{Professor: Jefferson}
\date{}

\begin{document}

\maketitle
\vspace{-1cm}

\begin{multicols}{2}

\section*{Parte 1: Compreensão}

\textbf{1. Qual é a profissão inicial de Lacie Pound?}  
\textcolor{red}{Ela trabalha em um escritório (funcionária administrativa).}

\vspace{0.2cm}

\textbf{2. Qual sistema de avaliação é utilizado?}  
\textcolor{red}{Um sistema de pontuação social baseado em avaliações de 1 a 5 estrelas.}

\vspace{0.2cm}

\textbf{3. Qual evento Lacie deseja participar?}  
\textcolor{red}{O casamento de Naomi, sua amiga de infância.}

\vspace{0.2cm}

\textbf{4. O que acontece quando a pontuação cai?}  
\textcolor{red}{A pessoa perde privilégios, acesso a serviços e status social.}

\vspace{0.2cm}

\textbf{5. Impacto no transporte e habitação?}  
\textcolor{red}{A pontuação define acesso a melhores moradias e meios de transporte.}

\vspace{0.2cm}

\textbf{6. Como o comportamento é influenciado?}  
\textcolor{red}{As pessoas agem de forma artificial para agradar e evitar avaliações negativas.}

\vspace{0.2cm}

\textbf{7. Por que o nome "Nosedive"?}  
\textcolor{red}{Refere-se à queda rápida da pontuação e da vida social da protagonista.}

\vspace{0.2cm}

\textbf{8. Objeto que representa a avaliação?}  
\textcolor{red}{O celular, onde as avaliações são feitas.}

\vspace{0.2cm}

\textbf{9. Relação Lacie e Naomi?}  
\textcolor{red}{Mostra superficialidade, baseada em status e aparência, não em amizade verdadeira.}

\vspace{0.2cm}

\textbf{10. Final do episódio?}  
\textcolor{red}{Lacie perde tudo e, na prisão, finalmente se expressa livremente, simbolizando libertação.}

\section*{Parte 2: Análise Crítica}

\textbf{11. Como o sistema afeta o comportamento?}  
\textcolor{red}{As pessoas vivem para agradar, perdendo autenticidade e espontaneidade.}

\vspace{0.2cm}

\textbf{12. Comparação com redes sociais?}  
\textcolor{red}{Assim como curtidas e seguidores, há busca por aprovação e construção de imagem idealizada.}

\vspace{0.2cm}

\textbf{13. Vantagens e desvantagens?}  
\textcolor{red}{Vantagem: incentivo a boas interações. Desvantagem: pressão social e perda de liberdade.}

\vspace{0.2cm}

\textbf{14. Sociedade valoriza aparência?}  
\textcolor{red}{Sim, muitas vezes aparência e status são mais valorizados que autenticidade.}

\vspace{0.2cm}

\textbf{15. Personagem mais autêntico?}  
\textcolor{red}{A caminhoneira (Susan), pois ignora o sistema e vive com liberdade.}

\section*{Parte 3: Conexão Pessoal}

\textbf{16. Estratégia para manter avaliação?}  
\textcolor{red}{Equilibrar educação e autenticidade, sem viver apenas para agradar.}

\vspace{0.2cm}

\textbf{17. Pressão nas redes sociais?}  
\textcolor{red}{Resposta pessoal, mas geralmente envolve mudança de comportamento para aceitação.}

\vspace{0.2cm}

\textbf{18. Influência do episódio?}  
\textcolor{red}{Leva à reflexão sobre uso consciente das redes sociais.}

\vspace{0.2cm}

\textbf{19. O que mudaria?}  
\textcolor{red}{Resposta pessoal; pode incluir um final mais crítico ou maior impacto social.}

\end{multicols}

\end{document}
